\cotefile{101}
\scanpage{30}
% L1–2
On peut \unsure{remplacer|expliciter|expliquer}, dans le cas général, le th. de dualité locale par une suite spectrale.

% L3–5
\uline{Lemme 6} $A'\longrightarrow\deleted{B}A$ un hom. d'anneaux, $N$ un $\deleted{B}A$-module, $M$ un $A'$-module. On a une suite spectrale cohomologique
% L6
\[
\operatorname{Ext}_{A'}^{\unsure{*|n|\bullet}}(N,M)\Longleftarrow
\operatorname{Ext}_{\deleted{B}A}^{p}(N,\operatorname{Ext}_{A'}^{q}(\deleted{B}A,M))
\]
% margin note written sideways
\margin{\uline{Cor} 1 Soit $x\in A$ non diviseur de $0$ dans $A$, ni dans $M$. Alors on a $\operatorname{Ext}_{A}^{p}(N,M)\simeq\operatorname{Ext}_{A/x}^{p-1}(N,M/xM)$ pour tt $N$, \uline{Cor 2} \unsure{Avec} une $M$-\unsure{suite}.}

% L7–8
Elle provient de la formule d'associativité
\[
\operatorname{Hom}_{A'}(N,M)=\operatorname{Hom}_{\unsure{B|A}}(N,\operatorname{Hom}_{A'}(B,M))
\]
% L9–10
\diagram{\begin{tikzcd}[ampersand replacement=\&]
{\unsure{\mathcal C|\mathscr C|\mathcal E}}^{\unsure{A'|A}}\arrow[r,dash,"{h_{\deleted{B}A}^{A'}}"]\arrow[rr,bend right=23,swap,"{h_N^{A'}}"]
 \&\mathcal C^{\deleted{B}A}\arrow[r,"{h_N^{\deleted{B}A}}"] \&{\unsure{\mathcal G|\mathscr G}}^{\deleted{B}A}
\end{tikzcd}}
$h_{\deleted{B}A}^{A'}$ transforme injectifs en injectifs

\centerline{- - - - -}

% L11–13
\uline{Prop. 12} \inserted{Soit $A$ un anneau quotient d'un Coh. Mac. $A'$ muni d'un mod. fdt $\Omega_{A'}$.} Il y a une suite spectrale canonique de type cohomologique
% L14
\[
\boxed{(E^{\deleted{\unsure{*}}i}(M))=(H^{\deleted{\unsure{d}}-i}(M))\underset{\unsure{p}}{\Longleftarrow}
 E_2^{pq}(M)=D(\operatorname{Ext}_{A}^{p}(M,\Omega_{A}^{(q)}))}
\]
% slanted
\margin{\uline{\unsure{Module fondamental|Modules fondamentaux}} \uline{pour un} \uline{anneau} de \uline{\unsure{polynomes}}\unsure{....}}

% L15–16
On applique le lemme 6 avec $N$\deleted{\unsure{$=$}} remplacé par $M$, et $M$ remplacé par $\Omega_{A'}$.

% archive sheet number: 29
\scanpage{31}
% L1–2
6. \uline{Caractérisation des modules fondamentaux d'anneaux Coh. Mac.} \uline{Anneaux de Gorenstein}.

% L3–4
\deleted{\uline{Théorème 5} Soient $A$ un anneau local noeth, $M$ un $A$-module}

% L5–6, large cross over the first version
\deleted{\uline{Prop. 13} Soient $A$ un anneau local noeth\inserted{ de corps rés. $k$,} $M$ un $A$-module, $n$ un entier $>0$, on suppose}
% L7
\[
\deleted{\operatorname{Ext}^{n-1}(k,M)=0}
\]
% L8–11
\deleted{Alors les conditions suivantes sont équivalentes}

\deleted{a) $\operatorname{Ext}^{n}(k,M)\simeq k$ et $\operatorname{Ext}^{n+1}(k,M)=0$}

\deleted{b) $N\longrightarrow\operatorname{Ext}^{n}(N,M)$ est dualisant}

\deleted{c) $H^{n}(M)$ est un module dualisant}

% L12–13
\deleted{On a alors $n=\dim A$, les idéaux pr. associés à $A$ sont \unsure{tous de dim $n$|tous de la dim.|étant de dim. $n$}, et $\lg_{\mathfrak p}(M)=\lg_{\mathfrak p}(A)$ pour tt tel $\mathfrak p$}

% L14–19, moved clauses connected by lines
\deleted{\uline{Cor. 1} Soit $\mathfrak p$ un idéal pr. minimal de $A$, les conditions précédentes équivalent alors} \deleted{à la conjonction des} \deleted{aux conditions : $n=\dim A$,} \inserted{\deleted{Supposons $n=\dim A$.}}

\deleted{$0<\lg_{\mathfrak p}M\leq\lg_{\mathfrak p}A$ et à une \unsure{conditions}} \deleted{$\unsure{\operatorname{Ext}^{n-1}(k,M)\ne0|\operatorname{Ext}^{n-1}(k,M)=0}$, $n=\dim A$,} \inserted{\deleted{et $0<\lg_{\mathfrak p}M\leq\lg_{\mathfrak p}A$}}

\deleted{et à l'une des conditions équivalentes suivantes}

% L20–22
\deleted{a') $\operatorname{Ext}^{n+1}(k,M)=0$}

\deleted{b') $N\longrightarrow\operatorname{Ext}^{n}(N,M)$ est exact}

\deleted{c') $H^{n}(M)$ est injectif} $\downarrow$

% L23–26, fresh version at foot
\deleted{\uline{Lemme 7}} \inserted{\uline{Proposition 13}} Soit $A$ anneau local\inserted{ de dim $n$}, $M$ un $A$-module, \deleted{\illegible{1}}

\deleted{\illegible{2}} Supposons $\boxed{\operatorname{Ext}^{n-1}_{A}(k,M)=0}$, \inserted{on a $\boxed{\operatorname{Ext}^{n}(N,M)\simeq\operatorname{Hom}(N,H^{n}(M))}$ pour $N$ de long. finie} \deleted{Alors} et les conditions suivantes sont équivalentes

% L27–29
a') $\operatorname{Ext}^{n+1}(k,M)=0$

b') $N\longrightarrow\operatorname{Ext}^{n}(N,M)$ est exact ($N$ de long finie)

c') $H^{n}(M)$ est injectif

% archive sheet number: 30
\scanpage{32}
% L1–5
\deleted{\unsure{Réciproquement|Prouvons|Démonstration} :}

$a\Rightarrow b\Rightarrow c$ trivial,

$c\Rightarrow a$ car $\unsure{N\mapsto|\mathstrut N\simeq}\operatorname{Ext}^{n}(N,M)$ sera exact, donc $N\longrightarrow\operatorname{Ext}^{n+1}(N,M)$ sera exact à gauche, donc $\operatorname{Ext}^{n+1}(k,M)\to H^{n+1}(M)$ sera injectif, or $H^{n+1}(M)=0$.

% L6–10
\uline{Corollaire 1} Sous ces conditions, \inserted{on a $H^n(M)\simeq D(A)^m$ pour certain $m$,} \deleted{les} idéaux premiers \deleted{et \unsure{ce} sont les idéaux pr. associés \unsure{$M$},} \deleted{de} \inserted{à $0$ dans} $A$ sont tous de dim $n$, et on a pour tout \deleted{\unsure{idéal}} tel $\mathfrak p$
\[
 \lg_{\mathfrak p}(M)=m\lg_{\mathfrak p}(A)
 \qquad \begin{aligned}\operatorname{Ext}^{n}(k,M)&\simeq k^{m}\\
 \operatorname{Ext}^{i}(k,M)&=0\ \text{si } i>n\end{aligned}
\]

% L11–12
\deleted{\unsure{Corollaire|Cor.} 2 \unsure{Si} $0<\lg_{\mathfrak p}M\leq\lg_{\mathfrak p}A$, les conditions a) b) c) \unsure{qui précèdent} équivalent à b') $\unsure{\operatorname{Ext}^{n}(N,M)|\operatorname{Ext}^{n}(M,M)}$ \unsure{est} c') $H^{n}(M)$ est dualisant.}

% margin note written sideways
\margin{\uline{Cor 2} Les conditions précédentes jointes à $0<\lg_{\mathfrak p}(M)\leq\lg_{\mathfrak p}(A)$ équivalent aux conditions \deleted{\unsure{équiv.|équivalentes}} suivantes compte tenu de $\operatorname{Ext}^{n-1}(k,M)=0$.\newline a) $\operatorname{Ext}^{n}(k,M)=k$, $\operatorname{Ext}^{n+1}(k,M)=0$\newline b) $N\longrightarrow\operatorname{Ext}^{n}(N,M)$ est dualisant\newline c) $H^{n}(M)$ est dualisant}

% L13–19
\uline{Remarque} \deleted{On ignore} \inserted{[\unsure{à reporter|reporter|injectif} à la fin du N° \unsure{+}]} Il semble que ces conditions impliquent que \deleted{\unsure{$B$}} $A$ est Coh.-Mac, et \deleted{\unsure{$B=\operatorname{Hom}$}} $M\simeq\Omega^{m}$, où $\Omega$ est un module fondamental pour $A$, \inserted{[et que l'hyp. $\operatorname{Ext}^{n-1}(k,M)=0$ soit $=$ redondante dans b')]} Pour le prouver, procédons par réc. sur $n$, il suffirait de prouver que pour $n>1$, l'annulateur de $\mathfrak m$ dans $M$ est nul. Voici un résultat partiel dans ce sens

% margin note written sideways
\margin{\uline{Cor 3} Les conditions du Cor 2 sont entraînées par l'une quelcon. des cond. suivantes\newline d) $\operatorname{Ext}^{n}(k,M)\simeq k$, $M$ fidèle\newline \unsure{d')bis)|d bis)} $\operatorname{Ext}^{\unsure{r|n}}(\unsure{A/\mathfrak n|A/\mathfrak m|A/\mathfrak p|A/\mathfrak N},M)$ est \deleted{\unsure{de|co}} \unsure{comonogène|monogène}, $M$ fidèle\newline \uline{Prouver que} \deleted{\unsure{$M$ soit}} ces idéaux pr. associés à $M$ sont tous de dim $n$. \uline{[\unsure{question}]}}

% L20–24
\uline{Proposition 14} \deleted{Sous les conditions précédentes,} \inserted{Supposons $m=1$, \unsure{ou} $A$ un quotient de régulier.} supposons de plus $M$. Coh. Mac. Alors les conditions a) b) c) \deleted{\unsure{sont encore équivalentes|sont chacune équivalente|sont devenues équivalentes}} à :

$A$ est Coh. Mac. et $M\simeq\Omega_A^{m}$, où $\Omega_A$ est un module fondamental pour $A$

% L25–26
C'est trivial si $n=0$, procédons par réc. sur $n$. Voir les notes jaunes - - - - -

% archive sheet number: 31
\scanpage{33}
% L1–4
\deleted{\uline{Proposition 15}} \deleted{\unsure{Supp. les}}

\uline{Cor. 1} \deleted{Sous ces \unsure{condit.} les conditions équivalentes des} \inserted{Supposons $M$ Coh. Mac sur $A$ de dim $n$.} Alors les conditions \deleted{\unsure{préci|précédentes}} suivantes sont équivalentes \inserted{à : $A$ est Coh. Mac et $M\simeq\Omega_A$ :}

% L5–8
a) b) c) du cor. 2 à prop 13 ; d / d)bis) du cor. 3

\deleted{d) $\operatorname{Ext}^{n}(k,M)=k$, $M$ est fidèle}

\deleted{e) $\operatorname{Ext}^{n}(\unsure{A/\mathfrak n|A/\mathfrak m|A/\mathfrak p|A/\mathfrak N},M)$ est monogène, $M$ est fidèle}

\deleted{f) Il suffit de prouver que}

% L9–12
\deleted{\uline{\unsure{Cor 2}} Soit $M$ comme dans le cor. 1, \unsure{Posons} $x_1\ldots x_n$ un syst. de paramètres de $M$. Alors \unsure{en|on} les conditions précédentes équivalent à la suivante}

\inserted{Soit $(x_1\ldots x_n)$ un syst. de paramètres de $M$.}

% L13
f) $M/x_1M\unsure{\ldots|+\cdots} + x_nM$ \deleted{est} est comonogène.

% L14–19
\uline{Remarque}. Il est possible que dans le cor. 1, \deleted{\unsure{la}} on puisse remplacer \deleted{Ext} la cond. : $M$ Coh. Mac. par la condition : $\operatorname{Ext}^{n-1}(k,M)=0$. Il semble $=$ que cette condition soit surabondante dans d bis) \deleted{(\unsure{donc aussi en d)}) et} [et à fortiori dans \unsure{b)|b), c)} et d)] \deleted{\illegible{1}} et \deleted{\unsure{en|$=$}} dans f.

% L20–22
La proposition suivante est de nature plus élémentaire

\uline{Proposition 15} Soit $A$ un anneau loc noeth. Les conditions suivantes sont équiv.

% L23–26
(a) $A$ est un module fondamental pour $A$.

(b) $H^n(A)$ est dualisant

(c) ${\unsure{H|h}}^n(A)$ est injectif

(d) $H^n(A)$ est comonogène et $A$ est \unsure{équidim|d'épaisseur|équidimensionnel}

% archive sheet number: 32
\scanpage{34}
% L1–5
(i) signifie que $D(A)\simeq H^n(A)$, i.e. (ii).

ii) $\Rightarrow$ (iii), d'autre part (iii) $\Rightarrow H^n(A)\simeq D(A)^m$ i.e. $A^m$ est \deleted{dualisant} fondamental. En vertu de prop. 9, on a $\lg_{\mathfrak p}(A^m)=\lg_{\mathfrak p}(A)$, d'où $m=1$, d'où (iii) $\Rightarrow$ (i). \deleted{Il reste à \unsure{prouver}}

% L5–8
Evident (i) $\Rightarrow$ (iv) par prop. 9, réciproque : $H^n(A)$ \unsure{sera|s'in.|comon.} $\subset D(A)$, l'annulateur de $H^n(A)$ est nul d'après l'hyp., donc $H^n(A)=D(A)$.

% L9–10
\deleted{\uline{Définition 2} Un tel $A$ est dit anneau de Gorenstein. S'il est de plus Cohen-Macaulay}

% L11–15, crossed draft
\deleted{\uline{Th. 5} Un anneau de Gorenstein est Coh. Mac}

\deleted{En effet, il est équidim, \unsure{donc} si \unsure{sa dim $n$}} \deleted{Réc. sur $n=\dim A$, trivial si $n=0$. Si $n>0$, $A$ étant équidim, il y a un $x\in A$ \unsure{non|t non} diviseur de $0$. \unsure{Considérons} Alors \deleted{la suite exacte} \inserted{la suite spectrale}}

% L16–20
\deleted{(où $N$ est un \deleted{\unsure{$A$}} \inserted{$A/x$}-$A/xA$-\unsure{module}) \unsure{de} \unsure{donne}}
\[
 \deleted{0\longrightarrow H^{n}(}\qquad \deleted{H^{n-1}(\,)}
\]
\[
 \deleted{\operatorname{Ext}_{A}^{*}(N,A\deleted{\unsure{/x}})\Longleftarrow
 \operatorname{Ext}_{A/x}^{*}(N,\operatorname{Ext}_{A}^{*}(A/x,A))}
\]
\deleted{\unsure{donne}, puisque} d'après le cor. au lemme 6,

% L21–25
\uline{Prop 16} bis \inserted{$+\,A$ est Cohen-Macaulay}, ces conditions équivalent aussi aux suivantes

\deleted{\unsure{b')|a')} \unsure{(ii)bis} Le foncteur $N\longrightarrow\operatorname{Ext}_{A}^{n}(N,A)$ est dualisant} \deleted{\unsure{ou que}} \deleted{[ou simplement exact]} \inserted{\deleted{exact}}

\inserted{c')} \deleted{\unsure{d)}} $\operatorname{Ext}^{n}(k,A)\simeq k$

\deleted{\unsure{e)}} $\operatorname{Ext}^{n+1}(k,A)=0$

% archive sheet number: 33
\scanpage{35}
% L1–4
\[
\left.\begin{aligned}
 \text{b')}\quad &\operatorname{Ext}^{n}(N,A)\ \text{dualisant en }N\\
 \text{c')}\quad &\operatorname{Ext}^{n}(\unsure{M|N}\unsure{,}A)\ \text{exact en }N
\end{aligned}\right\}
\quad\begin{gathered}\text{ici il est facile de}\\\text{voir que ça implique}\\\text{déjà Coh. Mac.}\end{gathered}
\]
c'') $\operatorname{Ext}^{n+1}(M,A)=0$, et $\operatorname{Ext}^{n-1}(k,A)=0$

e) $\operatorname{Ext}^{n}(k,A)=k$ et $\deleted{\unsure{\operatorname{Ext}^{n+1}(k,A)=0}}$ $M$ est Coh. Mac

% L5–7
\deleted{f) $\unsure{A/(x_1\ldots x_n|A/(x_1,\ldots,x_n)}$}

\deleted{Immédiat sans rien \unsure{spécifier} en début du N° (où $A$ sera équidimensionnel) \unsure{N°, 7, 9}}

% L8–10
\uline{Corollaire 1} \unsure{Ça équivaut|S'équivalent|Ce qui équivaut} aussi à :

e') $\operatorname{Ext}^{n}(A/\unsure{\mathfrak n|\mathfrak N},A)$ est comonogène

e'') $A/(x_1\ldots x_n)A$ est comonogène [i.e autodual]

% L11–12
\uline{Corollaire 2} Supposons $A$ quotient d'un régulier $A'$. Alors ces conditions équivalent : $\operatorname{Ext}_{A'}^{n'-n}(A,A')\simeq A$

% L13–15
\uline{Remarques} \deleted{\unsure{Dans}} Il est possible que dans \unsure{tous|toutes} ces conditions, l'hyp. Cohen-Macaulay soit surabondante.

% L16–18
\uline{Déf 2} Anneau de Gorenstein $=$ Anneau Coh. Mac. $A$ tel que $A$ \unsure{soit} fond.

Ex : les \uline{réguliers}

% L19–23
\uline{Proposition 17} (i) Un localisé d'un Gorenstein (supposé quotient d'un régulier) est Gorenstein

(ii) Kif Kif pour un quotient par l'idéal engendré par une partie d'un syst. des paramètres.

% L24–25
\uline{Cor} Quotient d'un régulier par un idéal engendré par une \inserted{partie d'un} syst. des paramètres.

% archive sheet number: 34
\scanpage{36}
% L1–5
$M$ module de type fini sur l'anneau local $A$ de dim $n$

$k$ corps résiduel de $A$

$\unsure{\mathfrak n|\mathfrak N}$ idéal de définition fixé

$x_1\ldots x_n$ syst. de paramètres fixé

$\mathfrak p$ un idéal premier de $A$ de corang $n$, fixé \qquad $N$ variable de long. finie, $N_0\ne0$ fixé

% L6–22, two columns connected by implication arrows
\deleted{\unsure{(iii)|(iv)}}
\diagram{\begin{tikzcd}[ampersand replacement=\&,row sep=.75em,column sep=.3em,font=\small]
\text{(i)}\arrow[d,Rightarrow]\arrow[dddd,bend right=40,Rightarrow] \&
\parbox[t]{12em}{\raggedright $\operatorname{Ext}_A^{n-1}(k,M)=0$ et l'une des conditions \inserted{a) à d)} suivantes [qui sont équivalentes moyennant la première, car on aura $\operatorname{Ext}^{n}(N,M)\simeq\operatorname{Hom}(N,H^n(M))$ etc ]} \& \parbox[t]{18em}{\raggedright a) $\operatorname{Ext}^{n+1}(k,M)=0\qquad \operatorname{Ext}^{n}(k,M)\simeq k$\newline
\deleted{\unsure{b)}} $*$b) $N\to\operatorname{Ext}^{n}(N,M)$ est dualisant\newline
$*$c) $N\to\operatorname{Ext}^{n}(N,M)$ est exact, $0<\lg_{\mathfrak p}M\leq\lg_{\mathfrak p}A$\newline
b') $H^n(M)$ dualisant\newline
c') $H^n(M)$ injectif, $0<\lg_{\mathfrak p}M\leq\lg_{\mathfrak p}A$\newline
d) $\operatorname{Ext}^{n+1}(k,M)=0$, $0<\lg_{\mathfrak p}M\leq\lg_{\mathfrak p}A$}\\
\text{(i bis)}\arrow[u,dotted,Rightarrow,"{\substack{\text{[si }M\,\deleted{\unsure{\text{Coh. Mac.}}}\\\text{\inserted{équidim.}}\,\text{]}}}" description]\arrow[d,Rightarrow] \&
\parbox[t]{12em}{\raggedright $\operatorname{Ext}_{A}^{n-1}(k,M)=0$ et l'une des \deleted{deux cond.} équiv. suivantes} \& \parbox[t]{18em}{\raggedright 
$*$e) $\operatorname{Ext}^{n}(k,M)\simeq k$, $M$ fidèle\newline
$*$f) $\operatorname{Ext}^{n}(A/\mathfrak n,M)$ est comonogène, $M$ fidèle}\\
\text{(ii)}\arrow[d,Rightarrow]\&
\parbox[t]{12em}{\raggedright Sous l'une des formes suivantes, qui sont équivalentes} \& \parbox[t]{18em}{\raggedright 
a) $N\to\operatorname{Ext}^{n}(N,M)$ est dualisant\newline
b) $N\to\operatorname{Ext}^{n}(N,M)$ exact, $0<\lg_{\mathfrak p}M\leq\lg_{\mathfrak p}A$}\\
\text{(iii)}\& \operatorname{Ext}^n(k,M)\simeq k,\quad M\ \text{fidèle}\\
\text{(iv)}\& \operatorname{Ext}^n(A/\mathfrak n,M)\simeq D(A/\mathfrak n),\quad M\ \text{fidèle}
\end{tikzcd}}
% L18–22, right of (iii)–(iv)
\margin{\deleted{\illegible{3}} NB, \unsure{même|Ici} si \inserted{$A$ et} $M$ sont Coh. Mac. il ne suffit pas ici de \unsure{considérer} \deleted{\illegible{4}} la \unsure{longueur}, car $\operatorname{Ext}^n(\unsure{M|N},M)=\operatorname{Hom}(\unsure{M|N},H^n(M))$ en dim $0$ \unsure{p.ex.}, si on ajoute des \unsure{élémts} nilpotents à $A$ \deleted{\illegible{2}}}

% L23–24
peut-être suffit-il même que $\operatorname{Ext}^n(A/\mathfrak n,M)$ soit comonogène et $M$ de support $\operatorname{Spec}(A)$

% L25–26
(v) L'annulateur de $\mathfrak m$ dans $M/x_1M+\ldots+x_nM$ est de dim $1$, $M$ est fidèle

% L27–35
Ces conditions sont vraies si $\boxed{A\text{ est Coh Mac et }M\simeq\Omega_A}$

Pour prouver la réciproque, il suffirait de montrer que dans les cas \unsure{iv)|ii)|(i bis)} à iii) (si on peut \unsure{iii), i)bis!?|(i bis) faible}) que \inserted{si $n>0$,} l'annulateur de $\mathfrak m$ dans $M$ est nul, \deleted{dans le cas (iv), qu'il en est de même dans \unsure{aussi|avec} $M$ et $A$,} Dans (v), \deleted{\unsure{que} si $n>0$ et si} \inserted{peut songer à} commencer par prouver que si $n>0$, et la \uline{première} condition de (v) est vérifiée, alors l'annulateur de $\mathfrak m$ dans $M$ est nul [d'où déjà que $M$ est Coh-Mac par récurrence, puis que $\operatorname{Ext}^n(A/\underline{x},M)\simeq M/\underline{x}M$ et $\operatorname{Ext}^n(k,M)\simeq k$, d'où le fait que $H^n(M)$ est le dual d'un $A/J$, $J$ intersection d'idéaux premiers de corang $n$

% archive sheet number: 35
\scanpage{37}
% L1–6, continuation
d'où \inserted{(remplaçant $A$ par $A/J$)} \unsure{(i)|(ii)} $+$ le fait que $M$ est Coh. Mac. Or de ceci on déduit facilement que $A$ est Coh. Mac. et $M\simeq\Omega_A$. On peut essayer une voie analogue pour (iv), en y déduisant de la \uline{première} hyp. \inserted{$+$ la condition du support} que l'annulateur de $\mathfrak m$ dans $M$ \uline{et} $A$ \deleted{sont} nuls, d'où par récurrence que $M$ \uline{et} $A$ sont Coh. Mac, \deleted{\unsure{ce}} qui ramène comme précédemment à un cas du (i)

% L7–13
Démonstration des équivalences dans (i)

$b\Leftrightarrow b'\qquad c\unsure{\Rightarrow|\Leftrightarrow}c'$, \qquad $a\Rightarrow b$ \deleted{\unsure{$\Rightarrow$}} triviaux

$b\Rightarrow c$ car $b\Leftrightarrow b'$, donc $D(H^n(M))\simeq\widehat A$ et $\lg_{\mathfrak p}(A)=\lg_{\mathfrak p}(M)$

$c'\Rightarrow b'$ car $c'\simeq\quad D(H^n(M))\simeq\widehat A^{m}$ et $m\lg_{\mathfrak p}(A)=\lg_{\mathfrak p}(M)$ d'où $m=1$

\margin{\deleted{\unsure{$b\Leftrightarrow b'$|$b'$}} \unsure{$c\Rightarrow c'$} \unsure{équiv.}}

$c\Rightarrow d$, car $N\to\operatorname{Ext}^{n+1}(N,M)$ est exact à gauche, d'où $\operatorname{Ext}^{n+1}(k,M)\to H^{n+1}(M)$ \uline{inj}.

$d\Rightarrow c'$ trivial \deleted{\unsure{donc}} // d'où $a\Rightarrow b\Leftrightarrow b'\Leftrightarrow c\Leftrightarrow c'\Leftrightarrow d$

% L14–16
\[
\begin{aligned}
b&\Rightarrow\operatorname{Ext}^n(k,M)\simeq k,\quad
  \text{et b')}\Rightarrow H^n(M)\text{ est fidèle donc }M\text{ est fidèle, d'où}\\
 &\phantom{\Rightarrow}\operatorname{Ext}^n(A/\unsure{\mathfrak n|\mathfrak N},M)\simeq D(A/\mathfrak n)
\end{aligned}
\]
$(b,b')\Rightarrow e$ et $f$, donc $(b,b',d)\Rightarrow a$

% L17–21
Reste à prouver que f) $\Rightarrow$ a). \inserted{Si $M$ est Coh. Mac.} Or
\[
\begin{aligned}
\operatorname{Ext}^n(A/\mathfrak n,M)&=\operatorname{Hom}(A/\mathfrak n,M)\\
\operatorname{Ext}^n(k,M)&=\operatorname{Hom}(k,\unsure{M|\mu})\\
\text{d'où}\qquad \operatorname{Ext}^n(k,M)&=\text{annulateur de }\mathfrak m\text{ dans }\operatorname{Ext}^n(A/\mathfrak n,M)
\end{aligned}
\]
donc f) $\Rightarrow$ e), et e) $\Rightarrow H^{n}(M)\subset D(A)$.

% L21–26, lines across earlier argument
\deleted{Donc \unsure{supposons}, si} \inserted{Or l'annulateur} \deleted{$J$ est l'annulateur de $H^n(M)$, on a $H^n(M)\simeq D(A')$ $A'=A/J$. Or $J$ est l'intersection des idéaux premiers de corang $n$ associés à $M$. Or $H^n(M)$ \deleted{\unsure{ne}} change pas si on remplace $M$ par \deleted{\illegible{1}} $M/M_1$, où $M_1$ est l'intersection des noyaux des $M\to M_{\mathfrak p}$ ($\mathfrak p$ de corang $n$).} de $H^n(M)$ est celui de $M$ ($M$ Coh. Mac. donc \uline{équidimensionnel}) donc $M$ étant fidèle, il est nul ; cela signifie (supposons $A$ complet, ce qui est loisible) que $H^n(M)=D(A)$, i.e. b').

% L27–31
\deleted{\unsure{Dans cette question|Donc}} Il reste donc essentiellement les

\uline{Problème} A Si \inserted{dans} $M/x_1M+\ldots+x_nM$, l'annulateur de $\mathfrak m$ est de dim $1$, et \inserted{$(\dim M=n)$} prouver que dans $M$ il est nul

\uline{Problème} B Si \inserted{dans} $\operatorname{Ext}^n(A/\mathfrak n,M)$, \unsure{le annulateur|l'annulateur} de $\mathfrak m$ est de dim $1$, et \inserted{$(\dim M=\dim A=n)$} prouver que dans $M$ \inserted{et $A$} il est nul

% archive sheet number: 36
\scanpage{38}
% L1–3
On en déduirait des \uuline{corollaires} : dans les deux cas $M$ est Coh-Mac, et si on désigne par $A'$ son anneau d'homothéties, $A'$ est Coh. Mac, et $M\simeq\Omega_{A'}$.

% L4–6
\uline{Problème} C Supposons \inserted{$\operatorname{Ext}^{n}(k,M)\ne0$ i.e. $H^n(M)\ne0$.} $\operatorname{Ext}^{n-1}(k,M)=\operatorname{Ext}^{n+1}(k,M)=0$, (la dernière \unsure{condition|condit.} \unsure{équivalant|équivalente} à : $H^n(M)$ injectif). \inserted{ou : $N\to\operatorname{Ext}^n(N,M)$ exact} Est-il vrai que $A$ est Coh. Mac. et $M\simeq\Omega_A^{m}$

% L7–10
[On notera qu'on a un hom naturel $M\to\Omega_A^{m}$, il faut montrer que c'est un isomorphisme]. \uline{Ici encore}, on est ramené (par le dévissage à la Serre) à prouver que l'annulateur de $\mathfrak m$ dans $M$ \uline{et} $A$ \deleted{est} sont nuls.

% archive sheet number: 37
\scanpage{39}
% L1–8
\uline{Théorème} On a les implications marquées sur le diagramme des notations. De plus, si on remplace dans (i) et (i bis) la relation $\operatorname{Ext}^{n-1}(k,M)=0$ par : $M$ est Coh. Mac. et si dans (v) on suppose de plus $M$ Coh. Mac., alors \deleted{toutes} les conditions \deleted{(i)} (i bis) (v) deviennent équivalentes à : $A$ est Coh. Mac. et $M\simeq\Omega_A$ (on peut alors oublier ii iii iv, qui pour $M$ Coh. Mac sont des rédites de (i) et (i bis)

% L9–11
$M$ est équidim, donc (i) $\Leftrightarrow$ (i bis). Prouvons par récurrence que alors $A$ est Coh. Mac. et $M\simeq\Omega_A$. Pour $n=0$ c'est trivial [$H^0(M)=M$ est dualisant].

% L12–22
Supposons $n>0$, et le th. prouvé pour $n'<n$. \deleted{\unsure{les}} Comme $H^n(M)$ est fidèle, \deleted{les idéaux pr. associés à $0$ de $A$ sont tous de dim. $n$, donc \unsure{et à}} \inserted{$M$ est fidèle, et étant} Coh. Mac, il s'ensuit que les idéaux pr. associés à $M$ et à $A$ sont les $=$s, et qu'ils sont tous de $=$ \deleted{dimension} \inserted{corang} $n$. Comme $n>0$, il y a un $x$ qui n'est contenu dans aucun d'eux, donc qui n'est pas div. de $0$ dans $A$ ou dans $M$. \deleted{\unsure{Donc on|D'après}} Ecrivons alors

% L23–26
\[
 \operatorname{Ext}_{A}^{\unsure{*|\bullet}}(N,M)\Longleftarrow
 \operatorname{Ext}_{A/x}^{*}(N,\operatorname{Ext}_{A}^{*}(A/x,M))
\]
On a $\operatorname{Ext}_{A}^{0}(A/x,M)=0$, $\operatorname{Ext}_{A}^{1}(A/x,M)=M/xM$, d'où
\[
 \operatorname{Ext}_{A}^{i}(N,M)\simeq\operatorname{Ext}_{A/x}^{\unsure{i-1|i1}}(N,M/xM)
\]
pour tout $N$. Par suite, la condition \unsure{(i a)|(i'a)}

% archive sheet number: 38
\scanpage{40}
% L1–5
\margin{\unsure{$\vert\,\vert$|11}}
reste vraie pour le module $M/xM$ sur $A/xA$. Or ce dernier est encore Coh. Mac, et l'hyp. de récurrence s'applique. Donc \deleted{$M/xM$} $A/xA$ est Coh. Mac, donc $A$ l'est. D'autre part $M/xM$ \deleted{$\simeq\Omega_{A/xA}$ donc} est fond. pour $A/xA$, d'où

% L6–8
\[
\begin{aligned}
 D(M/xM)&=\text{annulateur de }x\text{ dans }D(M)\\
 \rotatebox{90}{$\simeq$}&\\
 H^{n}(A/xA)&=\text{annulateur de }x\text{ dans }H^n(A).
\end{aligned}
\]

% L9–18
Donc l'hom naturel [\unsure{s'ass à|déduit de} $H^n(M)\simeq H^n(A)\otimes M=D(A)$]
\[
 H^n(A)\longrightarrow D(M)
\]
est tel qu'il induit un isom sur les éléments annulés par $x$, donc (lemme) il est injectif, ou par transposition
\[
 M\to\Omega_A
\]
est surjectif. D'autre part il induit un isom pour les $H^n$, i.e. le $H^n$ du noyau est $0$, donc ce noyau est de dim $<n$. Comme $M$ est équidim, ce noyau est nul.

% L19–24
\deleted{Prenons maintenant (v), en supposant $M$ Coh. Mac.}

\deleted{Procédons} \inserted{Supposons} \deleted{$M=$ \unsure{socle|Spec}$(M)$} \deleted{\inserted{pour}} $M/x_1M\unsure{+\ldots|\ldots}+x_nM$ a un annulateur de $\mathfrak m$ de dim $1$, $M$ \deleted{Coh. Mac.} \inserted{Coh. Mac \uline{fidèle}.}

\deleted{\unsure{Prouvons}} Alors \inserted{(\unsure{§ 1|cf 1|cf. § 1}, prop 4, cor. 1),} $\operatorname{Ext}^{n}(\deleted{\unsure{A/|A/\underline{x}}}\underline{k},M)=\deleted{\operatorname{Hom}(A/\underline{x},\unsure{H^{n}(M)|M^{n}(M)|M/\underline{x}M})}\inserted{\operatorname{Hom}(k,M/\underline xM)\simeq k}$, donc on est ramené à \unsure{(i* e)|(i e)|(i et e)}

% archive sheet number: 39
\scanpage{41}
% L1–6
$A$ anneau local, quotient d'un régulier $A'$. $n=\dim.\,\mathrm{Krull}\ A$.

$\underline{k}=A/\unsure{\mathfrak r|\mathfrak m}(A)$

$\unsure{\mathfrak n|\mathfrak N}$ un idéal de définition.

$(x_1,\ldots,x_n)$ un système de paramètres.

$\mathfrak p$ un idéal premier de dim $n$.

\uline{Conditions}

% L7–9
(i) $\operatorname{Ext}^{n}(\underline k,M)\simeq k,\ \deleted{[}\operatorname{Ext}^{n+1}(\underline k,M)=\deleted{[}\operatorname{Ext}^{n-1}(\underline k,M)=0\deleted{]]}$
% slanted
\margin{\unsure{on peut coupler|on peut compléter} comme on veut !}
\margin{$\operatorname{Ext}^{n}$ et $\operatorname{Ext}^{n+1}$ ne suffisent pas ni $\operatorname{Ext}^{n}$ et $\operatorname{Ext}^{n-1}$ ni $\operatorname{Ext}^{n-1}$ et $\operatorname{Ext}^{\unsure{n+1|n-1}}$. Contre-exples $A/\mathfrak p$, $A$ et $A/\mathfrak p$ rég. $A/J$, $\Omega_{A}$, où $A$ de dim $1$ non équ. $A^m$, où $A$ rég.}

(ii) $\operatorname{Ext}^{n}(\underline k,M)\simeq k$ et $M$ \uline{fidèle} [\deleted{\unsure{ou}} ou \unsure{seult|seulem$^{\mathrm t}$} de support tt]

% L10–13
\deleted{(iii) $\operatorname{Ext}^{n+1}(\underline k,M)=\operatorname{Ext}^{n-1}(\underline k,M)=0,\ 0<\lg_{\mathfrak p}(M)\leq\lg_{\mathfrak p}A$}

\deleted{(iii bis) $\operatorname{Ext}^{\unsure{n-1|n}}(\underline k,M)=0$ $\left.\begin{smallmatrix}M\text{ fidèle}\\ A\text{ équidim}\end{smallmatrix}\right\}$ $0<\lg_{\mathfrak p}(M)\leq\lg_{\mathfrak p}(A)$}

\deleted{(iv) \unsure{bis} $\operatorname{Ext}^{n-1}(\underline k,M)=0$ et $\operatorname{Ext}^{n}(A/\mathfrak n,M)\simeq A/\mathfrak n$}

% L14–15
(ii bis) $\operatorname{Ext}^{n}(A/\mathfrak n,M)\simeq\widehat{A/\mathfrak n}$ et $M$ \uline{fidèle} [faut-il $\operatorname{Ext}^{n-1}(k,M)=0$ ? et $A$ équidim ?]

\inserted{suffit il \unsure{même|$=$} que $\operatorname{Ext}^{n}(A/\mathfrak n,M)$ soit \uline{comonogène} ?}

% L16–20
(iii)
\diagram{\scriptsize\begin{tikzcd}[ampersand replacement=\&,column sep=small,row sep=.25em]
 {} \& {\deleted{M\text{ fidèle},\, A\text{ équidim}}}
 \arrow[d,bend left=40]\arrow[ddd,bend left=48] \& {\deleted{\text{\textcircled{a}}}\ \deleted{\unsure{\text{b}}}}\\
 {\left\{\begin{array}{l}\operatorname{Ext}^{n-1}(\underline k,M)=0,\\ \text{et l'une des cond. équivalentes}\end{array}\right.}
 \& \operatorname{Ext}^{n+1}(\underline k,M)=0 \& \text{\textcircled{a}}\\
 |[draw,rounded corners=1em]|{0<\lg_{\mathfrak p}(M)\leq\lg_{\mathfrak p}(A)} \arrow[ur,no head]
 \& H^n(M)\text{ dualisant} \& \text{\textcircled{b}}\\
 {} \& H^n(M)\text{ injectif} \& \text{\textcircled{c}}
\end{tikzcd}}
% slanted
\margin{$H^n(M)$ \unsure{tt seul|tout seul} ne donne rien, car on peut rajouter à $M$ des modules de dim. $<n$, sans changer $H^n(M)$. De $=$, $H^n(M)$ et $M$ équidim \unsure{ne nous disent rien de nouveau}, car on peut \unsure{trouver des exemples}, \unsure{prendre} $M'\subset M$ où $\dim M/M'<n-1$. Alors $\operatorname{Ext}^{n-1}(k,M)=0$ est pratiquement \unsure{$\Omega_{A/J}^{m}$|comme $\Omega_{A/J}^{m}$|$\Omega_{A/\mathfrak I}$} \illegible{2} $M=A/J$ \unsure{pour $A/J$} \illegible{3}.}

% L21–25
(iv) Le foncteur $N\mapsto\operatorname{Ext}^{n}(N,M)$ est dualisant

(iv bis) Le foncteur $N\mapsto\operatorname{Ext}^{n}(N,M)$ est exact, $0<\lg_{\mathfrak p}(M)\leq\lg_{\mathfrak p}(A)$

\hfill ($N$ de long. finie)

$\uparrow$

(v) L'annulateur de $\mathfrak m$ dans $M/\sum x_iM$ est de dim $1$, et $M$ est \uline{fidèle}

$\boxed{\text{(vi)}}$ $A$ est C-M., et $M$ est $\simeq\Omega_A$

% L26–29
\margin{\uline{Pb} $\left.\begin{array}{l} M\ne0\\ \operatorname{Ext}^{n-1}(\underline k,M)=\left\{\begin{array}{l}\operatorname{Ext}^{n+1}(\underline k,M)=0\\ \text{ou }H^n(M)\text{ injectif}\end{array}\right.\end{array}\right.$ $\Longrightarrow A$ est C.M. et $M\simeq\Omega_A^m$.}
% slanted
\margin{$M\ne0$, $\operatorname{Ext}^{n}(N,M)$ exact $\Longrightarrow$}
% slanted
\margin{\unsure{trouver une condition|trouvons une condition|mais une condition} \unsure{multiplicative|multiple} \illegible{2} \unsure{caractérisant|effectivement} \unsure{les multiples} de $\Omega_{\unsure{A}}$ \unsure{$A^m$}.}

% L30–36
\deleted{Le seul \unsure{cas connu|cas} où $\operatorname{Ext}^{n-1}(k,M)\ne0$, \inserted{$M\ne0$,} est celui où $A/J$ (l'anneau équidim. associé à $A$) est C.M. et $M\simeq\Omega_A^m\times\unsure{A/\mathfrak n^{m'}|A/J^{m'}}$ [cela implique $\operatorname{Ext}^{n}(k,M)=k$ \unsure{si|s.} $n=1$].}

\deleted{\deleted{\unsure{Le seul cas}} Le seul cas connu où $\operatorname{Ext}^{n}(k,M)\simeq k$ est celui où $A/J$ est Coh-Mac et où $M=\Omega_{A/\mathfrak n}$, où $\mathfrak n$ est un idéal de $A$ tel que $A/\mathfrak n$ soit C-M. / ou, pour $n=0$, un sous-module de $\widehat A$.}

\margin{cela implique-t-il $M$ \deleted{\unsure{quasi équidim|équidim}} ?}

% archive sheet number: 40
\scanpage{42}
% L1–4
\uline{Proposition} Supposons le foncteur $N\to\operatorname{Ext}^{i}(N,M)$ ($N$ de longueur finie) \uline{exact \unsure{à ga|à g.}} Alors on a
\[
 \operatorname{Ext}^{i}(N\deleted{\unsure{,}},M)\simeq\operatorname{Hom}(N,H^{i}(M))
\]
\margin{[\deleted{a-t-on} \deleted{\unsure{alors $i$|alors $i=$|donc} \illegible{1}} $\operatorname{Ext}^{i}(k,M)$ ou $\operatorname{Ext}^{i-1}(k,M)=0$ ??]}

% L5–11
Supposons le foncteur \uline{exact}, alors ou bien il est identiquement nul, [i.e. \deleted{$H^{i}(M)=0$, \unsure{ou c'est|ou encore}} $\operatorname{Ext}^{i}(k,M)=0$, cela implique prob.$^{\mathrm{t}}$ $i>n$ \uuline{ou} $i<\deleted{\unsure{d|co}}\dim.\,\mathrm{coh.}\ M$], ou bien $i=n$, $\widehat{H^n(M)}\simeq A^{m}$. Alors $M$ est \uline{fidèle} $\operatorname{Ext}^{n}(k,M)\simeq k^{m}$. De plus, cela implique $\operatorname{Ext}^{i}(k,M)=0$ pr $i>n$.
\margin{[et \unsure{cela|ceci} implique-t-il que $A$ est \uline{Coh. Mac.}, et $M\simeq\Omega_A^m$] \deleted{\unsure{fidélité|fidèle}}}

\par\noindent\shortstack{\rule{.4\linewidth}{.6pt}\\[-1ex]\rule{.4\linewidth}{.3pt}}

% L12–16
% slanted
\margin{$M$ un \unsure{mod.|module|néc.} de type fini}
Supposons le foncteur $N\mapsto\operatorname{Ext}^{\deleted{\unsure{i|i+1}}i-1}(N,M)$ \uline{exact} \uuline{à droite}. Je dis qu'on a \deleted{ou bien} $\operatorname{Ext}^{\deleted{\unsure{i+1|i}}\unsure{i-1|i}}(k,M)=0$, ou bien $\operatorname{Ext}^{\deleted{\unsure{i+1|i-1}}i}(k,M)=0\Longleftrightarrow H^{i}(M)=0$.

En effet, on aura $\operatorname{Ext}^{i}(N,M)\simeq\operatorname{Hom}(N,H^{i}(M))$.

% L17
\diagram{\begin{tikzcd}[ampersand replacement=\&]
0 \arrow[r] \& M \arrow[r] \& J \arrow[r] \& M' \arrow[r] \& 0
\end{tikzcd}}

% archive sheet number: 41
\scanpage{43}
% L1–6
(i) $\operatorname{Ext}^{n}(k,A)\simeq k$

\deleted{(\unsure{i bis|ii}) $\operatorname{Ext}^{\unsure{n}}(A/\unsure{\mathfrak n|\mathfrak N},A)\simeq\widehat{\unsure{A|A/\mathfrak n}}$}

\deleted{(ii) $\operatorname{Ext}^{n-1}(k,A)\simeq0$}

(iii) $\operatorname{Ext}^{n+1}(k,A)=0$

(i bis) $\operatorname{Ext}^{n}(A/\mathfrak n,A)\simeq\widehat{A/\mathfrak n}$

\margin{Prouver que chacune des hyp. (i) (iii') (i bis) (vi) implique l'\unsure{équidimensionnalité|équidemensionnalité}.}
\deleted{\unsure{la seule|la seconde} $\lg\operatorname{Ext}^{n}(A/\mathfrak n,A)=\lg(A/\mathfrak n)$}

% L7–13
(iv) Le foncteur $N\mapsto\operatorname{Ext}^{n}(N,\unsure{A|M})$ est dualisant

(iv bis) \rule{.30\linewidth}{.4pt} exact ($N$ de long. finie)

(v) \deleted{\unsure{est injectif|est précisément}} $H^{n}(A)$ dualisant

(v bis) $H^n(A)$ injectif

(v ter) $H^n(A)$ comonogène

(vi) $A/\sum x_iA$ est autodual

(vii) $\overline A\simeq\overline\Omega_A$

% L14–17
\inserted{[signif. $N\mapsto\operatorname{Ext}^{n}(N,M)$ est exact]}
\diagram{\scriptsize\begin{tikzcd}[ampersand replacement=\&,column sep=tiny,row sep=small]
 \boxed{\text{(vii)}\Leftrightarrow\text{(v)}\deleted{\unsure{\text{(v)}|\text{(vi)}}}\Leftrightarrow\text{(v bis)}\Leftrightarrow\text{(v ter)}}
 \arrow[d,Rightarrow] \&
 {\wedge\ \text{$A$ équidim.}\ \deleted{\unsure{\text{(v)}|\text{(iv)}}}} \&
 \boxed{\text{(iv bis)}\Leftrightarrow\text{(iv)}}\arrow[l,Rightarrow]\arrow[d,Rightarrow]\\
 \boxed{\text{(\unsure{iii'|iii})}}\& {}\& {\text{(i bis)}\quad\text{(i)}\quad\text{(vi)}}
\end{tikzcd}}
$A$ est C-M. et $\overline\Omega_A=\overline A$

% L18–23
Prouvons que (vii) implique que $A$ est C-M. \inserted{On peut supposer $A$ \uline{complet}} \inserted{(\unsure{Récurrence|Raisonnement|Rien} non trivial si $n=0$)} En effet, $A$ est équidim. puisque le module $\overline\Omega_A$ l'est. Soit $x\in A$ un diviseur de zéro, il n'est pas diviseur de zéro dans $\Omega_A$. Je dis \deleted{donc} que $\overline\Omega_{A/xA}\simeq\Omega_A/x\Omega_A$

% L24–29
(C'est là \deleted{Ext$_{\!A'}^{\unsure{b|r|*}}$} le résultat spécial à la récurrence). En effet,
\[
 \overline\Omega_{A/xA}=\deleted{\unsure{H^{n-1}}(A/xA)\,\unsure{\operatorname{Hom}|D}}\ \text{or }\operatorname{Ext}_{A'}^{\unsure{n-1|n+1|r-n}}(A/x,A')
\]
or
\[
 \deleted{\unsure{H^{n-1}}(A/xA)^{r-n+1}=\unsure{\operatorname{Hom}}}
\]
\[
 \operatorname{Ext}^{*}_{A'}(A/x,A')\Longleftarrow
 \operatorname{Ext}^{i}_{A}\bigl(A/x,\underbrace{\overset{\inserted{r-n}}{\operatorname{Ext}_{A'}^{j}(A,A')}}_{\Omega_A^{(j)}}\bigr)
\]
soit \deleted{d'où}

% L30–34
\uline{Lemme}
\[
 \Omega^{(*)}(A/x)\Longleftarrow\operatorname{Ext}^{i}_{A}(A/x,\Omega_A^{(j)})
\]
\margin{[Bien connu $=$ pour $A/x$ remplacé par un $A$-module quelconque]}
\deleted{ou \unsure{$A/x$|$Ax$}} \deleted{\unsure{est-il}}

d'où
\diagram{\begin{tikzcd}[ampersand replacement=\&,column sep=small]
 0\arrow[r] \& \operatorname{Ext}_{A}^{1}(A/x,\Omega_A)\arrow[r] \&
 \underbrace{\Omega_{A/x}}\arrow[r] \&
 \operatorname{Hom}^{\deleted{\unsure{0|q}}}_{A}(A/x,\Omega_A^{(1)})\arrow[r] \& 0 .
\end{tikzcd}}

% archive sheet number: 42
